\exercisesfor{2}

\begin{enumerate}
\def\labelenumi{\arabic{enumi}.}
\item
  采用定义 1 和第 6 号中的记号 T、J、\(T_{n}\)。令 t 为 T 的 p 阶齐次张量，且 \(\sigma\) 是 \(\{1, 2, \ldots, p\}\) 的一个排列。则 \(t - \sigma t \in J + T_{p-1}\)。（将其归结为 \(\sigma\) 是两个相邻整数的换位。）
\item
  设 \(\mathbf{K}\) 是一个域。令 \(\mathfrak{g}\) 为李 \(\mathbf{K}\) 代数，其 \(\mathbf{U}\) 包络代数。
\end{enumerate}

\begin{enumerate}
\def\labelenumi{(\alph{enumi})}
\item
  设 \(u, v\) 属于 U。若 \(u\) 的滤过是 \(> n\)，且 \(v\) 的滤过是 \(> p\)，则 \(uv\) 的滤过为 \(> n + p\)（使用定理 1）。
\item
  由 (a) 推出，U 的可逆元素只有标量。
\item
  由 (b) 推出，U 的 Jacobson 根为零。
\end{enumerate}

\begin{enumerate}
\def\labelenumi{\arabic{enumi}.}
\setcounter{enumi}{2}
\item
  设 K 是一个域。令 g 为李 K-代数，U 为其包络代数。 U 的左正则表示对应于 g 上的表示 \(\rho\)，作用在 U 空间上。证明集合 \(U_{+}\) 是 U 中不含常数项的元素所成的集合，在 \(\rho\) 下稳定，且 \(U_{+}\) 在 U 中没有关于 \(\rho\) 的补空间。特别地，\(\rho\) 不是半单的。（将在 § 6 的定理 2 中再次讨论。）
\item
  设 K 是一个域。
\end{enumerate}

\begin{enumerate}
\def\labelenumi{(\alph{enumi})}
\item
  验证存在一个具有基 \((x, y, z)\)，满足 \([x, y] = z\) , \([x, z] = [y, z] = 0\)。令 U 为 g 的包络代数。证明 U 的中心是由 1 和 z 生成的子代数。（对每个 \(\xi \in g\)，考虑将 ad 扩张为 g 的对称代数 S 上的导子，\(\xi\)，寻找被这些导子消去的 S 中元素；然后应用第 8 号末尾的注记。）
\item
  验证存在一个具有基 \((x, y, z)\)，满足 \([x, y] = y\) , \([x, z] = z\) , \([y, z] = 0\)。证明 g 的包络代数的中心退化为标量。（方法相同。）
\end{enumerate}

\begin{enumerate}
\def\labelenumi{\arabic{enumi}.}
\setcounter{enumi}{4}
\item
  设 \(\mathbf{K}\) 是特征为 \(p > 0\)（\(p\) 为素数）。令 \(\mathfrak{g}\) 为 \(\mathbf{K}\)，具有一组基，并规范地识别为其包络代数 U 的子模。对于 \(\phi\) 的 K-模 \(\mathfrak{g}\) 成为 \(p\)-映射的充要条件是 \(x \mapsto x^p - \phi(x)\) 是关于 \(\lambda \mapsto \lambda^p)\) 的半线性映射，它将 \(\mathfrak{g}\) 映入 U 的中心。推出，若 \((b_{\lambda})\) 是 \(\mathfrak{g}\)，则存在 \(p\) -映射定义在 \(\mathfrak{g}\) 上的充要条件是，对所有 \(\lambda\)，存在 \(c_{\lambda} \in \mathfrak{g}\) 满足 \((\operatorname{ad} b_{\lambda})^p = \operatorname{ad} c_{\lambda}\)；若如此，则存在唯一一个 \(p\) -映射 \(x \mapsto x^{[p]}\) 满足 \(b_{\lambda}^{[p]} = c_{\lambda}\) 对所有 \(\lambda\)。
\item
  设 g 是环 K 上的李 p-代数，满足 \(pK = \{0\}\)（p 为素数），U 为其包络代数，\(\sigma\) 是将 g 映入 U 的典范映射。令 J 为 U 的双边理想，由下列元素生成：\((\sigma(x))^{p} - \sigma(x^{[p]})\)，其中 x 遍历 g。结合代数 \(\tilde{U} = U/J\) 称为 g 的限制包络代数。映射 \(\sigma\) 在取商后定义映射 \(\tilde{\sigma}\)（称为典范映射）将 g 映入 \(\tilde{U}\) ，它是 p-同态（当 \(\tilde{U}\) 被视为李 p-代数时）。
\end{enumerate}

\begin{enumerate}
\def\labelenumi{(\alph{enumi})}
\item
  代数 \(\tilde{\mathbf{U}}\) 与映射 \(\tilde{\sigma}\) 是以下泛映射问题的解：对每个 \(p\) -同态 \(f\)，将 \(\mathfrak{g}\) 映入结合代数 \(\mathbf{B}\) 上的 \(\mathbf{K}\)（视为李 \(p\)-代数），存在唯一一个 K-同态 \(f'\)，将 \(\tilde{\mathbf{U}}\) 映入 \(\mathbf{B}\)（连同它们的结合代数结构）满足 \(f = f' \circ \tilde{\sigma}\)。
\item
  证明，若 \((b_{\lambda})_{\lambda \in \Lambda}\) 是 \(\mathfrak{g}\)（其中 \(\Lambda\) 完全有序），\(\tilde{\sigma}\) 是单射，且若 \(x\) 与 \(\tilde{\sigma}(x)\) 在 \(x\in \mathfrak{g}\) 中视为同一对象，则元素 \(\Pi_{\lambda}b_{\lambda}^{\nu_{\lambda}}\)（其中 \(\lambda\) 按递增顺序排列，\(\nu_{\lambda}\) 除有限多个外均为零，且 \(0\leqslant \nu_{\lambda} < p\) 对所有 \(\lambda\)）构成 \(\tilde{\mathbf{U}}\) 的一组基。（规范地将 \(\mathfrak{g}\) 与 \(\mathrm{U}\) 识别为子模，记 \(\phi (x) = x^{p} - x^{[p]}\) 对所有 \(x\in g\)。（对每个复合指标 \(\alpha = (\alpha_{\lambda})\in \mathbf{N}^{(\Lambda)}\)，令 \(\alpha_{\lambda} = \beta_{\lambda} + p\gamma_{\lambda}\) 满足 \(0\leqslant \beta_{\lambda} < p\)，并令
\end{enumerate}

\[
\mathrm{T} _ {\alpha} = (\Pi_ {\lambda} b _ {\lambda} ^ {\beta_ {\lambda}}) (\Pi_ {\lambda} (\phi (b _ {\lambda})) ^ {\gamma_ {\lambda}}).
\]

证明 \(T_{\alpha}\) 构成 U 的一组基；而满足 \(T_{\alpha}\) 的 \(\gamma = (\gamma_{\lambda}) \neq 0\) 构成 J 的一组基；注意 \(\phi(b_{\lambda})\) 属于 U 的中心。）

\begin{enumerate}
\def\labelenumi{\arabic{enumi}.}
\setcounter{enumi}{6}
\tightlist
\item
  设 g 是环 K 上的李 p-代数，满足 \(pK = \{0\}\)（p 为素数）。g 的导子 D 称为 p-导子，如果 \(\mathrm{D}(x^{p}) = (\mathrm{ad} x)^{p-1}\)。对所有 \(x \in g\)。每个内导子都是 p-导子。
\end{enumerate}

\begin{enumerate}
\def\labelenumi{(\alph{enumi})}
\item
  若 \(\mathbf{L}\) 是结合 K-代数，则 \(\mathbf{L}\) 的每个导子都是 \(p\)-导子，当 \(\mathbf{L}\) 被视为李 \(p\)-代数时。（使用第 1 节练习 19 的公式 (2)。）
\item
  设 g 有一组基。g 的导子为 p-导子的充要条件是它可延拓为 g 的限制包络代数的导子。推出，g 的 p-导子构成 g 的导子所成 p-代数的李 p-子代数。
\item
  若 D 是 \(p\) -导子，作用于 \(\mathfrak{g}\) ，则 D 的核是 \(p\) -子代数。\(\mathfrak{g}\) .
\item
  对于 g 的每个导子 D，\(\mathrm{D}(x^{p}) - (\mathrm{ad} x)^{p-1}\) . \(\bar{D}x\) 属于 g 的中心，对所有 \(x \in g\)（使用第 1 节练习 19 的公式 (2)，应用于 \(\mathrm{L} = \mathcal{L}(\mathfrak{g})\)）。
\end{enumerate}

\begin{enumerate}
\def\labelenumi{\arabic{enumi}.}
\setcounter{enumi}{7}
\item
  证明，若模 \(\mathfrak{g}\) 是单生成模的直和，定理 1 仍然成立。（在证明中用 g 的对称代数替换模 P。）
\item
  令\(k\)为含两个元素的域，\(\mathbf{V}\)为向量空间\(k^3\)，\((x_1, x_2, x_3)\)为其规范基，且\(\mathbf{K}\)为\(\mathbf{V}\)的外代数；它是一个定义在\(k\)上的 8 维交换代数。令\(\mathfrak{g}\)为具有基\((e_1, e_2, e_3, e_{12}, e_{13}, e_{23})\)的李 K-代数，使得
\end{enumerate}

\[
\left[ e _ {1}, e _ {2} \right] = \left[ e _ {2}, e _ {1} \right] = e _ {1 2} \quad \left[ e _ {1}, e _ {3} \right] = \left[ e _ {3}, e _ {1} \right] = e _ {1 3}, \quad \left[ e _ {2}, e _ {3} \right] = \left[ e _ {3}, e _ {2} \right] = e _ {2 3}
\]

且其他括号均为零。令 \(\mathfrak{h}\) 为 \(\mathfrak{g}\) 的理想，由 \(u = x_{1}e_{1} + x_{2}e_{2} + x_{3}e_{3}\) 生成。作为模，\(\mathfrak{h}\) 由 \(u\)、\([u, e_1] = x_2e_{12} + x_3e_{13}\) , \([u, e_2] = x_1e_{12} + x_3e_{23}\) , \([u, e_3] = x_1e_{13} + x_2e_{23}\)。令

\[
v = x _ {1} x _ {2} e _ {1 2} + x _ {1} x _ {3} e _ {1 3} + x _ {2} x _ {3} e _ {2 3}.
\]

\[
\phi \left(e _ {1}\right) = \phi \left(e _ {2}\right) = \phi \left(e _ {3}\right) = 0, \phi \left(e _ {1 2}\right) = x _ {3}, \phi \left(e _ {1 3}\right) = x _ {2}, \phi \left(e _ {2 3}\right) = x _ {1}.
\]

\begin{enumerate}
\def\labelenumi{(\alph{enumi})}
\setcounter{enumi}{1}
\item
  设 \(f\) 为从 \(\mathfrak{g}\) 到结合 K-代数的线性映射，满足 \(f([x,y]) = f(x)f(y) - f(y)f(x)\) 对所有 \(x, y\) 中的 \(\mathfrak{g}\)。证明 \(f(v) = f(u)^2\)。
\item
  由 (a) 和 (b) 推出，李代数 \(\mathfrak{g} / \mathfrak{h}\) 到其包络代数的典范映射不是单射。
\end{enumerate}

\begin{enumerate}
\def\labelenumi{\arabic{enumi}.}
\setcounter{enumi}{9}
\tightlist
\item
  设 \(\mathfrak{g}\) 是域上的有限维李代数，U 是其包络代数，\(\mathbf{U}_n\) 是 U 中滤过 \(\leqslant n\) 的元素所成的集合。
\end{enumerate}

\begin{enumerate}
\def\labelenumi{(\alph{enumi})}
\item
  设 \(x \in \mathrm{U}_n, y \in \mathrm{U}_n\) 是两个非零元素。证明存在 \(u \in \mathrm{U}, v \in \mathrm{U}\) 使得 \(ux = vy\)。（比较 \(\dim(\mathrm{U}_p x) = \dim \mathrm{U}_p\)，\(\dim(\mathrm{U}_p y) = \dim \mathrm{U}_p\) 和 \(\dim \mathrm{U}_{p+n}\)。推出 \(\mathrm{U}_p x \cap \mathrm{U}_p y \neq \{0\}\) 对于 \(p\) 足够大时成立。）
\item
  证明 U 存在一个左商域（《代数学》，第 I 章第 9 节练习 8），同时也是右商域。
\end{enumerate}

